\documentclass[11pt]{article}

\usepackage{cmap}
\usepackage[T2A]{fontenc}
\usepackage[utf8]{inputenc}
\usepackage[english,russian]{babel}

\usepackage{amsmath, amssymb, amsthm, mathtools}
\usepackage[margin=2.3cm]{geometry}
\usepackage{enumitem}
\usepackage{microtype}
\usepackage{titlesec}

\titleformat{\section}{\large\bfseries}{\thesection.}{0.6em}{}
\titlespacing*{\section}{0pt}{12pt}{5pt}

\setlength{\parindent}{0pt}
\setlength{\parskip}{3pt}

\tolerance=2000
\emergencystretch=2em

\theoremstyle{definition}
\newtheorem{problem}{Задача}

\theoremstyle{remark}
\newtheorem*{remark}{Замечание}

\DeclareMathOperator{\tr}{tr}
\DeclareMathOperator{\rank}{rank}
\DeclareMathOperator{\conv}{conv}
\DeclareMathOperator{\aff}{aff}
\DeclareMathOperator{\diag}{diag}
\DeclareMathOperator{\dom}{dom}
\newcommand{\D}{\mathrm{D}}
\newcommand{\R}{\mathbb{R}}
\newcommand{\Sym}{\mathbb{S}}
\newcommand{\norm}[1]{\left\lVert #1 \right\rVert}
\newcommand{\iprod}[2]{\langle #1,\, #2 \rangle}
\newcommand{\ones}{\mathbf{1}}

\begin{document}

\begin{center}
{\Large\bfseries Домашнее задание \#1}\\[3pt]
{\normalsize Методы оптимизации, ФПМИ, осенний семестр 2026--2027}\\[6pt]
{\normalsize\itshape Дедлайн: 6.10.2026, 23:59}\\[3pt]
{\normalsize\itshape Максимальный балл: 50}
\end{center}

\vspace{0.3cm}

\section{Дифференцирование и матричное исчисление}

\begin{problem}[2 балла]
Пусть $A \in \Sym^{n}$, $\varepsilon > 0$ и $f(x) = \dfrac{x^{\top} A x}{x^{\top} x + \varepsilon}$,
$f \colon \R^{n} \to \R$. Найдите $\nabla f(x)$. Чему равны стационарные точки $f$ и как они связаны с собственными векторами
матрицы $A$?
\end{problem}

\begin{problem}[2 балла]
Пусть $F(X) = X^{k}$, $F \colon \R^{n \times n} \to \R^{n \times n}$, где $k \ge 2$ ---
фиксированное целое число.
\begin{enumerate}[label=(\alph*), nosep]
  \item Найдите $D F(X)[H]$.
  \item Найдите $\nabla\, \tr\bigl(X^{k}\bigr)$.% относительно скалярного произведения Фробениуса.
\end{enumerate}
\end{problem}

\begin{problem}[4 балла]
Для $X \in \Sym^{n}_{++}$ через $X^{1/2}$ обозначается единственная матрица из $\Sym^{n}_{++}$, квадрат которой равен $X$; в спектральном разложении $X = Q \Lambda Q^{\top}$ она задаётся формулой $X^{1/2} = Q \Lambda^{1/2} Q^{\top}$. Рассмотрим отображение $F(X) = X^{1/2}$, $F \colon \Sym^{n}_{++} \to \Sym^{n}_{++}$.
\begin{enumerate}[label=(\alph*), nosep]
  \item Докажите, что дифференциал $d\bigl(X^{1/2}\bigr)$ удовлетворяет уравнению Сильвестра
        \[
        X^{1/2}\, d\bigl(X^{1/2}\bigr) \;+\; d\bigl(X^{1/2}\bigr)\, X^{1/2} \;=\; dX .
        \]
  \item Найдите $\nabla\, \tr\bigl(X^{1/2}\bigr)$ на $\Sym^{n}_{++}$.
\end{enumerate}
\end{problem}

\begin{problem}[Метод сопряжённых состояний, 5 баллов]
Инженер задаёт параметры $\theta \in \mathbb R^p$ некоторой конструкции. Состояние системы $x(\theta)$ определяется неявно, из уравнений равновесия $A(\theta)\,x = b(\theta)$; а оптимизировать нужно скалярный показатель $L(\theta) = c^{\top} x(\theta)$. Прямой подсчёт градиента требует решать свою линейную систему на каждый из $p$ параметров. Метод сопряжённых состояний использует то, что сам вектор $\partial x/\partial\theta_i$ не нужен --- нужна лишь его свёртка с $c$, и потому весь градиент удаётся получить ценой одного дополнительного решения системы, независимо от $p$. Это тот же приём, который отличает обратный режим автоматического дифференцирования от прямого.

Пусть $\theta \in \R^{p}$, матрица $A(\theta) \in \R^{n \times n}$ обратима в окрестности точки
$\theta$, вектор $b(\theta) \in \R^{n}$, вектор $x(\theta)$ --- решение системы
$A(\theta)\,x = b(\theta)$, и $L(\theta) = c^{\top} x(\theta)$ при фиксированном $c \in \R^{n}$.
\begin{enumerate}[label=(\alph*), nosep]
  \item Выразите $\partial x / \partial \theta_i$, а затем $\nabla L(\theta)$, через
        $\partial A / \partial \theta_i$ и $\partial b / \partial \theta_i$.
  \item Выпишите ответ явно для $A(\theta) = A_0 + \sum_{i=1}^{p} e^{\theta_i} A_i$ и
        $b(\theta) = b_0 + \sum_{i=1}^{p} \theta_i^{2} b_i$, где
        $A_0, \dots, A_p \in \R^{n\times n}$ и $b_0, \dots, b_p \in \R^{n}$ фиксированы.
  \item Опишите наиболее эффективный способ вычисления $\nabla L(\theta)$. Сколько линейных систем
        нужно при этом решить, то есть сколько раз посчитать произведение обратной матрицы на
        вектор?
\end{enumerate}
\end{problem}

\section{Выпуклые множества и оболочки}

\begin{problem}[Экспоненциальный конус, 4 балла]
Выражение $1 - e^{-u}$ --- стандартная модель насыщения: если события идут пуассоновским потоком с интенсивностью, пропорциональной вложенному усилию $u$, то это вероятность хотя бы одного успеха.
Так устроены кривые охвата в медиапланировании, вероятность обнаружения в задачах поиска и вероятность выявления дефекта при тестировании. Чтобы такие задачи попадали в область применимости методов внутренней точки, экспоненту оформляют как принадлежность одному стандартному конусу:
\[
K_{\exp} \;=\; \operatorname{cl}\Bigl\{ (u, v, w) \in \R^{3} \;:\; v > 0,\ \ v\, e^{u/v} \le w \Bigr\} .
\]
\begin{enumerate}[label=(\alph*), nosep]
  \item Докажите, что множество $K_{\exp}$ выпукло и является острым конусом.
  \item Покажите выпуклость множества $\{x \in \R^{n} : e^{\,a^{\top} x + b} \le c^{\top} x + d\}$.
\end{enumerate}
\end{problem}

\begin{problem}[3 балла]
Найдите $\conv\{x \in \{0,1\}^{n} : \ones^{\top} x = k\}$ и докажите ответ.
\end{problem}

\begin{problem}[5 баллов]
Обозначим через $O(n) = \{X \in \R^{n\times n} : X^{\top} X = I\}$ множество ортогональных матриц.
Найдите $\conv O(n)$ и докажите ответ.
\end{problem}

\section{Отделимость, грани и фасеты}

\begin{problem}[2 балла]
Для каждой пары множеств постройте отделяющую гиперплоскость явно или докажите, что строго
отделить их нельзя.
\begin{enumerate}[label=(\alph*), nosep]
  \item $C = \{x \in \R^{2} : \norm{x}_2 \le 1\}$ и точка $x_0 = (2, 0)$;
  \item $C = \{x \in \R^{2} : x \ge 0\}$ и $D = \{x \in \R^{2} : x_1 + x_2 \le -1\}$;
  \item $C = \{x \in \R^{2} : x_2 \ge e^{x_1}\}$ и $D = \{x \in \R^{2} : x_2 \le 0\}$.
\end{enumerate}
\end{problem}

\begin{problem}[2 балла]
Пусть $X_0 \in \Sym^{n}$ и $X_0 \not\succeq 0$. Постройте гиперплоскость, отделяющую $X_0$ от
конуса $\Sym^{n}_{+}$, явно предъявив её нормаль относительно скалярного произведения Фробениуса.
Можно ли отделить строго?
\end{problem}

\begin{problem}[Пермутоэдр и мажорирование, 7 баллов]
В задачах справедливого распределения ресурс делится между $n$ участниками, и естественно требовать, чтобы распределение было <<не менее неравномерным>>, чем некоторое наперёд заданное распределение $a$. Формализация этого --- отношение мажорирования, и оказывается, что множество всех допустимых распределений есть выпуклая оболочка перестановок исходного вектора. У этого многогранника экспоненциально много и вершин, и фасет одновременно --- случай, в котором ни одно из двух представлений не является экономным.

\emph{Матрицей перестановки} называется матрица $P \in \R^{n \times n}$, для которой $P_{ij} \in \{0, 1\}$, $P \ones = \ones$ и $P^{\top} \ones = \ones$. Пусть $a \in \R^{n}$ имеет попарно различные координаты, $a_{[1]} > a_{[2]} > \dots > a_{[n]}$ --- они же, упорядоченные по убыванию, и $P(a) = \conv\{ P_\sigma a : P_\sigma \ \text{--- матрица перестановки}\}$.
\begin{enumerate}[label=(\alph*), nosep]
  \item Найдите $\dim P(a)$ и $\aff P(a)$. Сколько у $P(a)$ вершин?
  \item Докажите, что для любого непустого собственного подмножества $S \subset \{1, \dots, n\}$
        неравенство $\sum_{i \in S} x_i \le \sum_{j=1}^{|S|} a_{[j]}$ допустимо для $P(a)$.
  \item Докажите, что каждое такое неравенство задаёт фасету $P(a)$. Сколько их всего?
\end{enumerate}
\end{problem}

\section{Выпуклые функции}

\begin{problem}[3 балла]
Пусть $h \colon \R^{n} \to \R$ выпукла и неотрицательна, а $g \colon \R^{n} \to \R$ вогнута и положительна.
\begin{enumerate}[label=(\alph*), nosep]
  \item Докажите, что функция $f(x) = h(x)^{2} / g(x)$ выпукла на $\dom h \cap \dom g$.
  \item Выведите отсюда выпуклость функции
        $\dfrac{\norm{Ax - b}_2^{2}}{1 - x^{\top} x}$ на открытом единичном шаре.
  \item Останется ли утверждение (a) верным, если убрать требование неотрицательности $h$?
\end{enumerate}
\end{problem}

\begin{problem}[2 балла]
Пусть $\alpha \in \R^{n}$, $\alpha \ge 0$ и $\sum_{k=1}^{n} \alpha_k \le 1$. Докажите, что функция
$f(x) = \prod_{k=1}^{n} x_k^{\alpha_k}$ вогнута на $\R^{n}_{++}$.
\end{problem}

\begin{problem}[Логарифмический барьер на конусе второго порядка, 3 балла]
Ограничение $x \in K$ можно перенести в целевую функцию, добавив к ней индикаторную функцию $\delta_K$. Такое слагаемое выпукло, но принимает бесконечные значения и нигде не гладко, поэтому численно работать с ним нельзя. Существуют методы, которые заменяют $\delta_K$ на гладкую функцию, конечную внутри $K$ и уходящую в $+\infty$ при приближении к границе; такое слагаемое называют \emph{барьером}. Чтобы вспомогательные задачи оставались выпуклыми, барьер обязан быть выпуклым.
Для конуса второго порядка
$Q = \{(x,t) \in \R^{n} \times \R : \norm{x}_2 \le t\}$ им служит
\[
f(x, t) \;=\; -\log\bigl(t^{2} - x^{\top} x\bigr),
\qquad \dom f = \{(x,t) : t > \norm{x}_2\} .
\]
Докажите выпуклость функции $f$.
\end{problem}

\begin{problem}[Симметричные функции спектра, 6 баллов]
Многие важные функции матриц зависят только от спектра: след, наибольшее собственное значение, сумма $k$ старших собственных значений, логарифм определителя, ядерная норма. Все они получаются подстановкой вектора собственных значений в симметричную функцию вектора. Оказывается, выпуклость при такой подстановке переносится автоматически, и считать гессиан по матричному аргументу не приходится вовсе.

Матрица $S \in \R^{n \times n}$ называется \emph{дважды стохастической}, если $S_{ij} \ge 0$,
$S \ones = \ones$ и $S^{\top} \ones = \ones$. Известно, что всякая дважды стохастическая матрица есть выпуклая комбинация матриц перестановок.
Функция $f \colon \R^{n} \to \R$ называется \emph{симметричной}, если $f(Px) = f(x)$ для любой
матрицы перестановки $P$. Через $\lambda(X) = (\lambda_1(X), \dots, \lambda_n(X))$ обозначается
вектор собственных значений матрицы $X \in \Sym^{n}$.
\begin{enumerate}[label=(\alph*), nosep]
  \item Докажите, что если $f$ выпукла и симметрична, а $S$ дважды стохастическая, то
        $f(Sx) \le f(x)$.
  \item Пусть $X = Q \Lambda Q^{\top}$ --- спектральное разложение с ортогональной $Q$. Покажите,
        что матрица $S$ с элементами $S_{ij} = Q_{ij}^{2}$ дважды стохастическая и что
        $\diag(X) = S\,\lambda(X)$.
  \item Выведите равенство
        $f(\lambda(X)) = \sup_{V} f\bigl(\diag(V^{\top} X V)\bigr)$, где супремум берётся по всем
        ортогональным $V$, и заключите отсюда, что $X \mapsto f(\lambda(X))$ выпукла на $\Sym^{n}$.
  \item Примените результат к $f(x) = \log \sum_k e^{x_k}$ и получите выпуклость
        $X \mapsto \log \tr e^{X}$.
\end{enumerate}
\end{problem}


\end{document}
